\documentclass[12pt,oneside]{book}

% ============================================================
% METADATOS
% ============================================================
\newcommand{\universidad}{Universidad de Buenos Aires (UBA)}
\newcommand{\facultad}{Facultad de Derecho}
\newcommand{\carrera}{Orientación Notarial y Registral}
\newcommand{\materia}{Derechos Reales}
\newcommand{\titulo}{Apuntes de Derechos Reales: El Derecho Real de Servidumbre}
\newcommand{\autor}{Isaac Edgar Camacho Ocampo}
\newcommand{\anio}{2026}

% ============================================================
% PAQUETES
% ============================================================
\usepackage[utf8]{inputenc}
\usepackage[T1]{fontenc}
\usepackage[spanish,es-nodecimaldot]{babel}
\usepackage[
    top=2.5cm,
    bottom=2.5cm,
    left=3cm,
    right=2.5cm,
    headheight=15pt
]{geometry}
\usepackage{amsmath}
\usepackage{amssymb}
\usepackage{mathtools}
\usepackage{graphicx}
\usepackage{float}
\usepackage{caption}
\usepackage{subcaption}
\usepackage{booktabs}
\usepackage{array}
\usepackage{longtable}
\usepackage{tabularx}
\usepackage{enumitem}
\usepackage{xcolor}
\usepackage[most]{tcolorbox}
\usepackage{fancyhdr}
\usepackage{ragged2e}
\usepackage[
    colorlinks=true,
    linkcolor=blue!50!black,
    citecolor=green!40!black,
    urlcolor=blue!60!black,
    pdftitle={\titulo},
    pdfauthor={\autor},
    pdfsubject={Apuntes universitarios},
    bookmarks=true
]{hyperref}

% ============================================================
% CONFIGURACIÓN
% ============================================================
\graphicspath{
    {./img/}
    {./graficos/}
    {./figuras/}
}
\setcounter{tocdepth}{2}
\setlength{\parindent}{0pt}
\setlength{\parskip}{0.5em}
\setlist[itemize]{
    topsep=0.3em,
    itemsep=0.2em
}
\setlist[enumerate]{
    topsep=0.3em,
    itemsep=0.2em
}
\clubpenalty=10000
\widowpenalty=10000

% ============================================================
% COMANDOS
% ============================================================
\newcommand{\abs}[1]{\left|#1\right|}
\newcommand{\norm}[1]{\left\lVert#1\right\rVert}
\newcommand{\vect}[1]{\mathbf{#1}}
\newcommand{\deriv}[2]{\frac{d#1}{d#2}}
\newcommand{\pderiv}[2]{\frac{\partial#1}{\partial#2}}

% ============================================================
% ENTORNOS / CAJAS
% ============================================================
\newtcolorbox{definition}[1]{
    enhanced,
    breakable,
    title={Definición: #1},
    fonttitle=\bfseries,
    colback=blue!3,
    colframe=blue!50!black,
    boxrule=0.8pt,
    arc=2mm
}

\newtcolorbox{important}{
    enhanced,
    breakable,
    title={Importante},
    fonttitle=\bfseries,
    colback=yellow!8,
    colframe=orange!70!black,
    boxrule=0.8pt,
    arc=2mm
}

\newtcolorbox{example}[1]{
    enhanced,
    breakable,
    title={Ejemplo: #1},
    fonttitle=\bfseries,
    colback=green!4,
    colframe=green!40!black,
    boxrule=0.8pt,
    arc=2mm
}

\newtcolorbox{formula}[1]{
    enhanced,
    breakable,
    title={Fórmula / Regla: #1},
    fonttitle=\bfseries,
    colback=purple!4,
    colframe=purple!50!black,
    boxrule=0.8pt,
    arc=2mm
}

\newtcolorbox{exercise}[1]{
    enhanced,
    breakable,
    title={Ejercicio: #1},
    fonttitle=\bfseries,
    colback=gray!5,
    colframe=gray!60!black,
    boxrule=0.8pt,
    arc=2mm
}

\newtcolorbox{note}{
    enhanced,
    breakable,
    title={Nota},
    fonttitle=\bfseries,
    colback=white,
    colframe=gray!50,
    boxrule=0.6pt,
    arc=2mm
}

% ============================================================
% CONFIGURACIÓN FINAL (encabezados y pies)
% ============================================================
\pagestyle{fancy}
\fancyhf{}
\fancyhead[L]{\nouppercase{\leftmark}}
\fancyhead[R]{\materia}
\fancyfoot[C]{\thepage}
\renewcommand{\headrulewidth}{0.4pt}
\renewcommand{\footrulewidth}{0pt}
\fancypagestyle{plain}{
    \fancyhf{}
    \fancyfoot[C]{\thepage}
    \renewcommand{\headrulewidth}{0pt}
}

% ============================================================
% DOCUMENTO
% ============================================================
\begin{document}

% ------------------------------------------------------------
% PORTADA
% ------------------------------------------------------------
\begin{titlepage}

\thispagestyle{empty}

\begin{center}

\vspace*{1cm}

{\large \textsc{\universidad}}

\vspace{0.2cm}

{\large \textsc{\facultad}}

\vspace{0.2cm}

{\large \carrera}

\vspace{3cm}

{\Huge \textbf{\materia}}

\vspace{1cm}

{\LARGE El Derecho Real de Servidumbre}

\vspace{0.8cm}

{\small \textit{Estudiar con constancia es el camino para comprender con profundidad.}}

\vspace{1.2cm}

\rule{0.70\textwidth}{0.5pt}

\vspace{0.5cm}

{\large Apuntes de estudio}

\vspace{0.5cm}

\rule{0.70\textwidth}{0.5pt}

\vfill

{\large \autor}

\vspace{0.3cm}

{\large \anio}

\end{center}

\vfill

\begin{flushleft}
\textit{Este es un modesto aporte a los estudiantes de ``Derechos Reales'' no pretende sustituir el material oficial de la materia.}
\end{flushleft}

\end{titlepage}

% ------------------------------------------------------------
% ÍNDICE
% ------------------------------------------------------------
\tableofcontents

% ------------------------------------------------------------
% CONTENIDO
% ------------------------------------------------------------

% ============================================================
\chapter{Concepto y objeto de la servidumbre}
% ============================================================

\section{Introducción}

La \textbf{servidumbre} es un \textbf{derecho real de disfrute sobre cosa ajena} que \textbf{vincula dos fundos entre sí}. A partir de su constitución, el dominio del titular de uno de los fundos se \textbf{imperfecciona}: el dominio que era absoluto pasa a ser un \textbf{dominio imperfecto}, porque la constitución de la servidumbre \textbf{desmembra el dominio} del titular.

\begin{note}
Las notas originales se apoyan en un esquema ilustrativo que representa mediante tierras y fondos (campos o potreras) las situaciones analizadas; se lo utiliza solo como forma de presentar la figura, no porque la servidumbre se limite a fondos agrícolas. Dicho esquema no se encuentra disponible en el material original, por lo que a lo largo del texto se describen en palabras las situaciones que ilustraba.
\end{note}
% TODO: imagen del esquema de fondos (Fundo A / Fundo B, Ruta Nacional N.º 2) mencionada en las notas originales y no disponible

\section{Los fundos dominante y sirviente}

La servidumbre es un derecho real que vincula a dos fundos entre sí. En el esquema de referencia se identifican dos fundos: el \textbf{Fundo A}, que adquiere la calidad de \textbf{fundo dominante}, y el \textbf{Fundo B}, que adquiere la calidad de \textbf{fundo sirviente}.

Ambos fundos se vinculan a partir de la constitución del derecho real, que implica que el titular del Fundo B (fundo sirviente) deberá prestar a favor del otro fundo cierta \textbf{utilidad}. Esa utilidad puede consistir, por ejemplo, en pasar, transitar o sacar agua, e incluso puede ser de mero recreo. En todos los casos se traduce en una obligación de:

\begin{itemize}
    \item \textbf{dejar que} el titular del fundo dominante haga algo, o
    \item \textbf{abstenerse} de realizar determinada acción o actividad que el titular del dominio tendría si no existiese la servidumbre.
\end{itemize}

\begin{important}
En ningún caso puede tratarse de una \textbf{obligación de hacer}: siempre consiste en una abstención o en \emph{permitir que el otro haga}.
\end{important}

\section{Definición legal}

\begin{definition}{Servidumbre (art. 2162 CCCN)}
La servidumbre es el derecho real que se establece entre dos inmuebles y que concede al titular del inmueble dominante determinada utilidad sobre el inmueble del fondo sirviente, que es el inmueble ajeno.
\end{definition}

\section{Objeto de la servidumbre}

El objeto de la servidumbre puede ser \textbf{todo el terreno} o \textbf{una parte determinada} del terreno; por ejemplo, si se tratara de transitar, el paso puede ejercerse por todo el fundo o por un sector delimitado. Siempre se trata de \textbf{dos inmuebles con distinta titularidad}, que tienen por objeto una parte material determinada o todo el inmueble.

\section{Desmembramiento del dominio}

La constitución de la servidumbre desmembra el dominio del titular del fundo sirviente, que deja de ser un dominio perfecto para convertirse en un \textbf{dominio imperfecto}. Esa imperfección constituye la carga que grava al fundo sirviente.

\section{Relevancia profesional}

\begin{note}
El régimen de servidumbres se vincula con la práctica del operador jurídico especializado en Derechos Reales y Derecho Notarial, quien debe:
\begin{itemize}
    \item redactar escrituras públicas constitutivas de servidumbres (paso, acueducto, vistas), con cláusulas precisas sobre el fundo dominante, el fundo sirviente, la utilidad concedida y el precio;
    \item asesorar al cliente que adquiere un inmueble sobre las cargas reales que lo gravan: si existe una servidumbre real constituida, el adquirente la asume automáticamente;
    \item patrocinar reclamos judiciales de servidumbre forzosa de tránsito o acueducto cuando un fundo queda sin salida a la vía pública o sin acceso al agua, fijando la indemnización con el juez;
    \item identificar el momento de extinción de una servidumbre (falta de uso por 10 años, muerte del titular en las personales, desaparición de la utilidad) para cancelar la registración correspondiente.
\end{itemize}
\end{note}

\section*{Cuadro Cornell: concepto y objeto}
\addcontentsline{toc}{section}{Cuadro Cornell: concepto y objeto}

\begin{table}[H]
\centering
\small
\begin{tabularx}{\textwidth}{
    >{\raggedright\arraybackslash}p{0.26\textwidth}
    >{\raggedright\arraybackslash}X
}
\toprule
\textbf{Preguntas / palabras clave} &
\textbf{Notas / respuestas} \\
\midrule

\textbf{¿Qué es la servidumbre y qué establece el art. 2162?} &
Es el derecho real que se establece entre dos inmuebles y que concede al titular del inmueble dominante determinada utilidad sobre el inmueble sirviente (ajeno). A partir de su constitución, el dominio del titular del fundo sirviente se desmembra y pasa a ser un dominio imperfecto. \\

\textbf{¿Cuál es la diferencia entre fundo dominante y fundo sirviente?} &
El fundo dominante es el que recibe la utilidad (derecho de paso, agua, vista). El fundo sirviente es el que debe soportar o abstenerse. La relación siempre vincula dos inmuebles con distinta titularidad. \\

\textbf{¿Sobre qué recae la servidumbre y puede consistir en una obligación de hacer?} &
Recae sobre todo el inmueble sirviente o sobre una parte material determinada. Nunca puede consistir en una obligación de hacer: siempre es una abstención o un permitir que el otro haga. \\

\midrule

\multicolumn{2}{>{\raggedright\arraybackslash}X}{
\textbf{Resumen:}
La servidumbre es un derecho real de disfrute sobre cosa ajena que vincula dos inmuebles de distinta titularidad: el dominante, que recibe la utilidad, y el sirviente, que la presta. Su objeto puede ser todo el inmueble o una parte determinada, y desmembra el dominio del sirviente, que pasa a ser imperfecto.
} \\

\bottomrule
\end{tabularx}
\end{table}

% ============================================================
\chapter{Clasificación de las servidumbres}
% ============================================================

La servidumbre debe constituirse siempre mediante \textbf{escritura pública}, porque se trata de un derecho real constituido sobre inmuebles.

\section{Servidumbres positivas y negativas}

La denominación que se da al tipo de utilidad permite clasificar las servidumbres en positivas y negativas.

\begin{table}[H]
\centering
\renewcommand{\arraystretch}{1.3}
\begin{tabularx}{\textwidth}{
    >{\raggedright\arraybackslash}X
    >{\raggedright\arraybackslash}X
}
\toprule
\textbf{Positivas} & \textbf{Negativas} \\
\midrule
La carga consiste en \textbf{soportar} el ejercicio de una acción.
&
El titular del fundo sirviente debe \textbf{abstenerse} de realizar alguna facultad. \\[0.5em]
\emph{Ejemplo:} soportar que pasen por todo o por un sector determinado del fundo.
&
\emph{Ejemplo:} servidumbre de vista; el titular sirviente no puede construir y afectar la vista del fundo dominante. \\
\bottomrule
\end{tabularx}
\end{table}

\section{Servidumbres reales y personales}

Esta clasificación \textbf{no se refiere a la distinción entre derechos reales y derechos personales}, que son dos esferas diferentes de derechos. Estas servidumbres son siempre reales: no existen servidumbres personales como instituto dentro de los derechos reales. Se trata de un instituto propio de los derechos reales que vincula dos fundos.

\begin{itemize}
    \item \textbf{Servidumbre personal:} algunas servidumbres (siempre reales) son constituidas, por voluntad de las partes o por ley, a favor de una persona determinada, teniendo en cuenta quién es el titular del fundo dominante.
    \item \textbf{Servidumbre real:} se constituye teniendo en cuenta el fundo, con independencia de quién sea su titular.
\end{itemize}

\section{Servidumbres forzosas (de origen legal)}

En el caso de la servidumbre real pueden presentarse fondos encerrados que no tienen acceso a la vía pública, o fondos que carecen de agua. En estos casos la servidumbre es de \textbf{origen legal}, porque el orden público interviene con rigor y exige e impone la servidumbre.

El orden público no admite, por ser desfavorable para todos y para la comunidad, un fundo que carezca de toda utilidad. El ejemplo concreto es la servidumbre de paso o de tránsito.

\begin{example}{Fundo encerrado (servidumbre legal de tránsito)}
Se tienen dos fundos y todo lo marcado en el esquema como tierra. Uno de ellos no tiene salida a la vía pública (en el ejemplo, la Ruta Nacional N.º 2), de modo que sus titulares no podrían ingresar y, si lograran hacerlo por algún medio aéreo, tampoco podrían salir; no tendría sentido. Lo mismo ocurre con los fundos que carecen de agua.

La servidumbre de origen legal impone necesariamente al titular del fundo que tiene una salida más directa a la vía pública (en este caso, a la Ruta Nacional N.º 2) tolerar y permitir que el titular del fundo dominante pase.
\end{example}

\begin{important}
La mirada del orden público está puesta en la \textbf{utilidad del fundo}, no en la calidad personal del titular de dominio. Cualquiera sea el titular, aun cuando mute el titular registral por enajenación o transmisión \emph{mortis causa} a herederos, tendrá que poder pasar.
\end{important}

\section{Servidumbres personales convencionales}

Se analiza también una servidumbre de tránsito, pero personal. En el esquema, el inmueble A tiene salida a la vía pública (la Ruta Nacional N.º 2) por un camino lateral, un camino rural que bordea los campos y que, después de muchos kilómetros, sale a la ruta.

No puede imponerse al titular del inmueble B que permita que A pase por su fundo porque a A le resulte más cómodo o su acceso sea más directo. El \textbf{orden público no está comprometido} en ese caso: solo impone la servidumbre cuando no existe salida a la vía pública o no existe acceso al agua.

\begin{example}{Servidumbre personal convencional}
El titular de A nunca tuvo dificultades y le resultó suficiente salir por el camino existente, con su caballo o su vehículo, hacia la ruta. Cierto día decide cambiar su actividad o vende, y el adquirente quiere realizar una actividad que le impone un gran tránsito de camiones; ese camino le resulta poco práctico y además debe recorrer muchos kilómetros para llegar a la vía pública.

El adquirente pacta entonces con el titular del fundo sirviente la constitución de una servidumbre, proponiéndole un camino determinado por el cual pueda pasar con sus camiones. Hay libertad y consenso: no existe imposición legal, sino un origen exclusivamente convencional. Las partes pactan el precio y la duración.

Si no pudieran ponerse de acuerdo en la indemnización, A podría solicitarla judicialmente y el precio sería fijado por el juez. Nada de esto ocurre con la servidumbre personal cuando el orden público está satisfecho: la salida a la vía pública existe y la servidumbre es del interés de A y, si le cierra la ecuación económica, también del interés de B.
\end{example}
% TODO: contenido ambiguo en las notas originales: se afirma que ``no hay imposición legal'' y a la vez que, sin acuerdo sobre la indemnización, el juez fija el precio

\section{Cuadro comparativo: servidumbre real y personal}

\begin{table}[H]
\centering
\small
\renewcommand{\arraystretch}{1.3}
\begin{tabularx}{\textwidth}{
    >{\raggedright\arraybackslash}p{0.22\textwidth}
    >{\raggedright\arraybackslash}X
    >{\raggedright\arraybackslash}X
}
\toprule
\textbf{Aspecto} & \textbf{Real (forzosa)} & \textbf{Personal (convencional)} \\
\midrule
Origen & Legal; orden público & Convencional; voluntad de las partes \\
Requisito & Fundo encerrado o sin agua & Solo el interés de las partes \\
Camino & El más directo a la vía pública & El que acuerden las partes \\
Precio & Fijado por el juez si no acuerdan & Libremente pactado \\
Transmisión & Con el inmueble (muta el titular) & Se extingue al transmitir \\
Duración & Mientras subsiste la utilidad & Temporal; la acordada \\
Plazo máximo sin estipular & Indeterminado (mientras sea útil) & Vitalicia (persona humana) / 50 años (persona jurídica) \\
\bottomrule
\end{tabularx}
\end{table}
% TODO: el cuadro original equipara ``real'' con ``forzosa'' y ``personal'' con ``convencional''; en el desarrollo oral estas categorías se presentan como clasificaciones distintas

\section*{Cuadro Cornell: clasificación}
\addcontentsline{toc}{section}{Cuadro Cornell: clasificación}

\begin{table}[H]
\centering
\small
\begin{tabularx}{\textwidth}{
    >{\raggedright\arraybackslash}p{0.26\textwidth}
    >{\raggedright\arraybackslash}X
}
\toprule
\textbf{Preguntas / palabras clave} &
\textbf{Notas / respuestas} \\
\midrule

\textbf{¿Servidumbre positiva vs. negativa?} &
Positiva: el titular sirviente debe soportar el ejercicio de una acción (por ejemplo, que pasen por su campo). Negativa: debe abstenerse de una facultad que de otro modo podría ejercer (por ejemplo, no construir para no tapar la vista del dominante). \\

\textbf{¿Cuándo es <<real>> y cuándo es <<personal>> la servidumbre?} &
Real: se constituye teniendo en cuenta el fundo, con independencia de quién sea su titular, y se transmite con el inmueble. Personal: se constituye teniendo en cuenta la persona del titular dominante y se extingue al transmitirse el inmueble. Ambas son siempre servidumbres reales, en el sentido de que pertenecen a los derechos reales. \\

\textbf{¿Qué es una servidumbre forzosa y cuándo procede?} &
Es de origen legal; surge por mandato del orden público cuando un fundo carece de salida a la vía pública (fundo encerrado) o de acceso al agua. El titular sirviente no puede negarla; solo puede reclamar la indemnización, incluso judicialmente si no hay acuerdo. \\

\textbf{¿Qué pasa cuando las partes acuerdan libremente la servidumbre?} &
Es una servidumbre convencional: no hay imposición legal, las partes pactan el camino, el precio y la duración. El orden público no interviene cuando ya existe salida a la vía pública o acceso al agua. \\

\midrule

\multicolumn{2}{>{\raggedright\arraybackslash}X}{
\textbf{Resumen:}
Las servidumbres se clasifican en positivas (soportar) y negativas (abstenerse), y en reales (atienden al fundo y se transmiten con él) y personales (atienden al titular y se extinguen al transmitirse). Las forzosas o legales proceden por orden público ante fundos encerrados o sin agua; las convencionales nacen del acuerdo de partes con libertad de precio y plazo.
} \\

\bottomrule
\end{tabularx}
\end{table}

% ============================================================
\chapter{Constitución y legitimación}
% ============================================================

\section{Forma requerida}

La servidumbre debe constituirse siempre mediante \textbf{escritura pública}, porque se trata de un derecho real constituido sobre inmuebles. La escritura pública es el \textbf{título suficiente} de constitución de la servidumbre.

\section{Duración de la servidumbre personal}

Como la servidumbre personal es siempre convencional en su origen, tendrá una duración y será necesariamente \textbf{temporal}. Si no se determina el plazo en el instrumento constitutivo (la escritura pública), se considera vitalicia cuando el titular es una persona humana.

\begin{formula}{Plazo de la servidumbre personal (art. 2165 CCCN)}
Se presume que la servidumbre personal es vitalicia cuando no se ha determinado un plazo en el instrumento constitutivo. En el caso de persona jurídica, el plazo máximo es de cincuenta (50) años.
\end{formula}

En la servidumbre forzosa (la servidumbre legal: tránsito, fondos encerrados, acueducto, sacar agua), el titular del fundo sirviente tiene derecho a recibir la \textbf{indemnización} y a no sufrir perjuicios.

\begin{important}
En ningún caso la servidumbre personal puede ser constituida judicialmente.
\end{important}

\section{Legitimación: quién puede constituir la servidumbre}

Pueden constituir la servidumbre:

\begin{itemize}
    \item el \textbf{titular de dominio};
    \item en el condominio, \textbf{todos los condóminos} que totalicen el cien por ciento de la titularidad de la cosa.
\end{itemize}

\subsection{Nudo propietario y usufructo}

Si el titular del fundo B hubiese constituido un usufructo, \textbf{no} puede constituir además una servidumbre: el nudo propietario solo conserva el poder de disposición, y la servidumbre implica un menoscabo de ese derecho, que él ha desplazado a favor del usufructuario. Cuando existe usufructo, la facultad corresponde al \textbf{usufructuario}, y la servidumbre estará necesariamente vinculada al tiempo que dure el usufructo, que es siempre un derecho temporal.

\subsection{Uso, habitación y anticresis}

Si el nudo propietario hubiese constituido, en lugar de usufructo, un derecho de uso, de habitación o de anticresis, podrá constituir la servidumbre en alguna de estas condiciones:

\begin{itemize}
    \item con la autorización del habitador, del usuario o del acreedor anticresista, según corresponda; o
    \item constituyéndola de modo que comience a regir como condición necesaria cuando se hayan extinguido el uso, la habitación o la anticresis.
\end{itemize}

\section*{Cuadro Cornell: constitución y legitimación}
\addcontentsline{toc}{section}{Cuadro Cornell: constitución y legitimación}

\begin{table}[H]
\centering
\small
\begin{tabularx}{\textwidth}{
    >{\raggedright\arraybackslash}p{0.26\textwidth}
    >{\raggedright\arraybackslash}X
}
\toprule
\textbf{Preguntas / palabras clave} &
\textbf{Notas / respuestas} \\
\midrule

\textbf{¿Qué forma es necesaria para constituir una servidumbre?} &
Siempre escritura pública, por tratarse de un derecho real sobre inmuebles. Este instrumento es el título suficiente de constitución. \\

\textbf{Art. 2165: ¿qué plazo tiene la servidumbre personal?} &
Si no se fija plazo: vitalicia para persona humana; máximo de 50 años para persona jurídica. La servidumbre personal nunca puede ser constituida judicialmente. \\

\textbf{¿Quién está legitimado para constituirla?} &
El titular de dominio, o todos los condóminos conjuntamente. El nudo propietario no puede mientras exista usufructo (el usufructuario sí, por el tiempo del usufructo). Ante uso, habitación o anticresis, puede constituirla con autorización del titular de esos derechos o con efecto diferido a su extinción. \\

\midrule

\multicolumn{2}{>{\raggedright\arraybackslash}X}{
\textbf{Resumen:}
La constitución exige escritura pública y legitima principalmente al titular de dominio (o a todos los condóminos). El nudo propietario queda excluido durante el usufructo. La servidumbre personal, sin plazo, es vitalicia (persona humana) o de hasta 50 años (persona jurídica), y nunca puede constituirse judicialmente.
} \\

\bottomrule
\end{tabularx}
\end{table}

% ============================================================
\chapter{Transmisión, derechos y obligaciones}
% ============================================================

\section{Transmisión de las servidumbres reales}

Las servidumbres reales, que son las que vinculan los fundos entre sí, se transmiten \textbf{conjuntamente con el inmueble}, ya enajene el inmueble el titular del fundo dominante o el titular del fundo sirviente. Como los vinculados son los fundos, quienquiera que sea el titular de uno u otro, la servidumbre se transmite del uno al otro con independencia de la mutación del titular.

\begin{formula}{Transmisión de las servidumbres (art. 2172 CCCN)}
Las servidumbres reales se transmiten conjuntamente con el inmueble al que acceden, tanto la constituida a su favor como la que lo grava. La servidumbre personal se extingue con la transmisión del inmueble.
\end{formula}

Cuando se enajena el inmueble, los titulares lo adquieren con la servidumbre. La personal, en cambio, se extingue con la transmisión, porque tanto en su constitución como en su duración se tuvo en cuenta a una persona determinada, el titular de A.

\section{Derechos y obligaciones del titular del fundo dominante}

\begin{formula}{Fundo dominante (arts. 2173 a 2175 CCCN)}
El titular del fundo dominante puede realizar en el inmueble sirviente todas las obras necesarias para el ejercicio y conservación de la servidumbre. Tiene a su cargo las mejoras necesarias, como mantener el camino en condiciones. Debe comunicar al titular del fundo sirviente cualquier perturbación, de hecho o de derecho, de que tome conocimiento, para que pueda ejercer la defensa.
\end{formula}

El titular del fundo dominante tiene la utilidad, pero debe realizar todas las mejoras necesarias; si se trata de un camino, debe repararlo y mantenerlo en condiciones. Tiene además la facultad de hacer en el inmueble sirviente todo aquello que sea necesario para que su derecho no se vea menoscabado; por ejemplo, si por un hecho natural (como la caída de árboles) se impide la salida, puede retirar esos árboles. Debe también comunicar al titular del fundo sirviente cualquier perturbación, de hecho o de derecho, para que este pueda ejercer la defensa.

\section{Derechos y obligaciones del titular del fundo sirviente}

\begin{formula}{Fundo sirviente (arts. 2180 y 2181 CCCN)}
El titular del fundo sirviente conserva la disposición jurídica y material del inmueble. Puede ejercer todas las facultades propias de su derecho de dominio, con el límite de no afectar el ejercicio de la servidumbre. Puede realizar obras, constituir otros derechos reales y mejoras, siempre que no impidan o menoscaben el derecho del titular del fundo dominante.
\end{formula}

Al titular del fundo sirviente, a quien en las servidumbres de origen legal o forzoso se le impone una carga como la de acueducto o tránsito, le quedan todas las facultades materiales y jurídicas propias del dominio. El \textbf{límite} es que no puede realizar nada que afecte el derecho del titular del fundo dominante.

La afección o imperfección de su dominio se manifiesta en que hay muchas cosas que no puede realizar: no puede constituir otros derechos reales ni derechos personales que de algún modo afecten el goce (la utilidad) que el titular del fundo dominante recibe a partir de la servidumbre. Puede, en cambio, realizar obras, constituir otros derechos reales y realizar mejoras, pero nunca hacer nada que afecte al titular del fundo dominante en el ejercicio de las facultades que le concede la servidumbre.
% TODO: en el cuadro Cornell original la prohibición de constituir derechos que afecten la utilidad figura entre las obligaciones del fundo dominante; en el desarrollo se la refiere al titular del fundo sirviente

\section{Inembargabilidad de la servidumbre}

\begin{formula}{Inembargabilidad (art. 2178 CCCN)}
La servidumbre no puede ser embargada ni vendida con independencia del inmueble al que accede, ya sea el inmueble dominante o el sirviente.
\end{formula}

A diferencia de otros derechos reales ya estudiados (las notas mencionan el caso de los frutos, cuyo derecho era ejecutable), el art. 2178 indica con claridad que en ningún caso la servidumbre en sí misma es ejecutable: no puede ser embargada ni vendida con independencia del fundo que grava ni del fundo que beneficia.
% TODO: referencia a ``los frutos'' ambigua en las notas originales

\section*{Cuadro Cornell: transmisión, derechos y obligaciones}
\addcontentsline{toc}{section}{Cuadro Cornell: transmisión, derechos y obligaciones}

\begin{table}[H]
\centering
\small
\begin{tabularx}{\textwidth}{
    >{\raggedright\arraybackslash}p{0.26\textwidth}
    >{\raggedright\arraybackslash}X
}
\toprule
\textbf{Preguntas / palabras clave} &
\textbf{Notas / respuestas} \\
\midrule

\textbf{Art. 2172: ¿cómo se transmiten las servidumbres?} &
Las reales se transmiten conjuntamente con el inmueble, tanto la constituida a favor como la que lo grava, con independencia de la mutación del titular. La personal se extingue con la transmisión del inmueble. \\

\textbf{Derechos y obligaciones del fundo dominante (arts. 2173--2175)} &
Puede realizar obras necesarias en el sirviente para el ejercicio de la servidumbre. Tiene a su cargo las mejoras necesarias (mantener el camino). Debe comunicar al sirviente toda perturbación de hecho o de derecho. \\

\textbf{Derechos y obligaciones del fundo sirviente (arts. 2180--2181)} &
Conserva todas las facultades materiales y jurídicas propias del dominio. Puede hacer obras y constituir otros derechos reales, con el límite de no afectar el ejercicio de la servidumbre. No puede constituir derechos que afecten la utilidad concedida. Tiene derecho a recibir indemnización, fijada judicialmente si no hay acuerdo. \\

\textbf{Art. 2178: ¿es ejecutable la servidumbre?} &
No. La servidumbre no puede ser embargada ni vendida con independencia del fundo, sea el dominante o el sirviente. Es inseparable del inmueble al que accede. \\

\midrule

\multicolumn{2}{>{\raggedright\arraybackslash}X}{
\textbf{Resumen:}
Las servidumbres reales se transmiten con el inmueble; la personal se extingue al transmitirse. El titular dominante puede hacer las obras necesarias, soporta las mejoras necesarias y debe comunicar perturbaciones; el sirviente conserva sus facultades de dominio con el límite de no afectar la servidumbre. La servidumbre es inembargable e inseparable del fundo.
} \\

\bottomrule
\end{tabularx}
\end{table}

% ============================================================
\chapter{Extinción de la servidumbre}
% ============================================================

\section{Causales de extinción}

Cuando se trata de fondos encerrados o sin agua, la servidumbre se mantiene mientras subsistan esas condiciones de encerramiento o de falta de agua.

\begin{important}
\textbf{Causales de extinción (art. 2182 CCCN).} Son causales de extinción de la servidumbre: la desaparición de toda utilidad para el fundo dominante; la falta de uso por diez (10) años; en las personales, la muerte del titular cuando no se ha fijado un plazo (o, cuando sea persona jurídica, el vencimiento del plazo máximo de cincuenta (50) años); y el vencimiento del plazo fijado en el instrumento constitutivo.
\end{important}

\begin{itemize}
    \item \textbf{Desaparición de la utilidad.} Si un fundo encerrado necesita salir a la vía pública y luego se construye un camino lateral, el interés del orden público desaparece y la servidumbre carece de objeto.
    \item \textbf{Falta de uso por diez años.}
    \item \textbf{Servidumbres personales.} Muerte del titular, cuando no se haya fijado plazo, ya que la servidumbre personal es siempre temporal y la vida humana es esencialmente temporal; en el caso de la persona jurídica rige el plazo máximo de cincuenta años cuando no se ha fijado un plazo especial.
\end{itemize}

\section{Efectos de la extinción}

Al extinguirse la servidumbre se extinguen todos los derechos personales o reales que pudieran estar vinculados con ella. El titular del fundo sirviente, que veía aminorado su derecho, lo restablece: pasa a tener nuevamente un \textbf{dominio pleno}, también llamado \textbf{dominio perfecto}, pues ha desaparecido la carga que desmembraba su derecho.

\section*{Cuadro Cornell: extinción}
\addcontentsline{toc}{section}{Cuadro Cornell: extinción}

\begin{table}[H]
\centering
\small
\begin{tabularx}{\textwidth}{
    >{\raggedright\arraybackslash}p{0.26\textwidth}
    >{\raggedright\arraybackslash}X
}
\toprule
\textbf{Preguntas / palabras clave} &
\textbf{Notas / respuestas} \\
\midrule

\textbf{Art. 2182: ¿cuáles son las causales de extinción?} &
1) Desaparición de toda utilidad para el dominante. 2) Falta de uso por 10 años. 3) En personales: muerte del titular (si no hay plazo) o vencimiento del plazo estipulado. 4) Persona jurídica: vencimiento del plazo máximo de 50 años. \\

\textbf{¿Qué efectos produce la extinción?} &
Se extinguen todos los derechos personales o reales vinculados. El dominio del titular sirviente se restablece como dominio pleno (perfecto): desaparece la carga que lo desmembraba. \\

\midrule

\multicolumn{2}{>{\raggedright\arraybackslash}X}{
\textbf{Resumen:}
La servidumbre se extingue por desaparición de la utilidad, falta de uso por 10 años, muerte del titular personal o vencimiento del plazo; al extinguirse, el dominio sirviente recupera su plenitud y caen los derechos vinculados.
} \\

\bottomrule
\end{tabularx}
\end{table}

% ============================================================
\chapter{Síntesis general}
% ============================================================

La servidumbre es un derecho real que siempre vincula dos fundos.

En ocasiones se tiene en cuenta la calidad del fundo, con independencia del titular, como en las servidumbres \textbf{forzosas o legales}, que deben ser constituidas por imperio de la ley, más allá de la voluntad del titular del fundo sirviente, porque el orden público ha puesto allí su mirada.

Por otro lado están las \textbf{personales}: son siempre servidumbres reales, pero se llaman personales porque en su constitución se tiene en cuenta la persona del titular. Por eso son convencionales: las partes acuerdan el monto de la indemnización, la manera de pasar por el inmueble y los derechos y obligaciones que surgen de una situación consensual que las vincula durante el tiempo que ambas hayan fijado para la duración del instituto.

\section*{Síntesis final}
\addcontentsline{toc}{section}{Síntesis final}

\begin{table}[H]
\centering
\small
\begin{tabularx}{\textwidth}{X}
\toprule
La servidumbre es un derecho real de disfrute sobre cosa ajena que vincula dos fundos entre sí: el dominante, que recibe la utilidad, y el sirviente, que la presta, desmembrando el dominio de este último y convirtiéndolo en un dominio imperfecto. Su objeto puede abarcar todo el inmueble sirviente o una parte material determinada. Se clasifica en positivas (el sirviente soporta un hacer) y negativas (debe abstenerse), y en reales (atienden al fundo con independencia del titular y se transmiten con el inmueble) y personales (atienden al titular y se extinguen al transmitirse). Las legales o forzosas (paso y acueducto) operan por mandato del orden público cuando un fundo carece de salida a la vía pública o de agua: el titular sirviente está obligado a soportarlas y solo puede reclamar indemnización judicial. Las convencionales nacen del acuerdo de partes con libertad de precio y plazo. La constitución siempre requiere escritura pública y legitima principalmente al titular de dominio; el nudo propietario está excluido durante el usufructo. La servidumbre es inembargable e inseparable del fundo (art. 2178). Se extingue por desaparición de la utilidad, falta de uso por 10 años, muerte del titular personal o vencimiento del plazo; y al extinguirse, el dominio sirviente recupera su plenitud. \\
\bottomrule
\end{tabularx}
\end{table}

\end{document}
